%==============================================================================
%  CHAPTER 11: Gradient Descent Variants
%  Part II: Optimization Theory
%==============================================================================
\documentclass[../../main.tex]{subfiles}

\begin{document}

\chapter{Gradient Descent Variants}
\label{ch:gradient-descent}

The learning rate is the single most important hyperparameter in deep learning. Set it too high, and your training explodes. Set it too low, and it never converges. This chapter teaches you to navigate that critical balance.

\begin{whycare}
Gradient descent is the heartbeat of neural network training. Every training run you'll ever do uses some variant of gradient descent. Mastering it means faster experiments, better models, and fewer 3 AM debugging sessions.
\end{whycare}

\begin{keyconcept}
Gradient descent\index{gradient descent|textbf}\index{batch gradient descent|see{gradient descent}} lies at the heart of modern machine learning optimization. While the basic algorithm is straightforward---repeatedly move in the direction opposite to the gradient---the practical challenge lies in choosing how to take these steps. This chapter explores the landscape of gradient descent variants, from classical gradient descent through momentum and Nesterov acceleration, and provides concrete comparisons that guide practitioner decisions.
\end{keyconcept}

\begin{realworld}
\textbf{Why This Matters:} Gradient descent is the engine of deep learning---every framework is built around it.

\textbf{Real Example:} PyTorch and TensorFlow are essentially sophisticated gradient descent engines with autodiff.

\textbf{Scale:} Large language models like GPT-4 are believed to require enormous computational resources, though exact figures are not publicly disclosed.
\end{realworld}

\begin{prerequisites}
\textbf{Before reading this chapter, you should be comfortable with:}
\begin{itemize}
    \item Stochastic gradient descent basics (Chapter~\ref{ch:stochastic-optimization})
    \item Gradient computation and the chain rule (see Book~1, Chapter~5)
    \item Exponential moving averages
\end{itemize}
\textbf{Review if needed:} Chapter~\ref{ch:stochastic-optimization} for SGD fundamentals
\end{prerequisites}

%------------------------------------------------------------------------------
\section{Introduction}
\label{sec:ch11-intro}

The previous chapter established stochastic gradient descent\index{stochastic gradient descent} (SGD)\index{SGD|see{stochastic gradient descent}} as the workhorse of large-scale optimization. We learned that replacing exact gradients with noisy estimates enables training on massive datasets while providing beneficial implicit regularization. However, vanilla SGD faces significant practical challenges:

\begin{itemize}
    \item \textbf{Slow convergence} in ill-conditioned loss landscapes where curvature varies dramatically across directions
    \item \textbf{Sensitivity to learning rate}: too large causes divergence, too small causes glacial progress
    \item \textbf{Difficulty navigating} saddle points and narrow ravines common in deep network loss surfaces
    \item \textbf{Per-parameter learning rates}: different parameters may require vastly different step sizes (addressed in Chapter~\ref{ch:advanced-optimizers})
\end{itemize}

This chapter focuses on the foundational momentum methods that address these challenges by accumulating gradient information over time to accelerate convergence and smooth the optimization trajectory.

To ground these ideas concretely, Section~\ref{sec:ch11-optimizer-showdown} compares gradient descent and momentum methods on the Rosenbrock function, illustrating how momentum improves convergence on ill-conditioned problems.

\begin{important}
\textbf{Chapter Organization:} This chapter covers gradient descent and momentum methods. For adaptive learning rate methods (AdaGrad, RMSprop, Adam, AdamW) and learning rate schedules, see Chapter~\ref{ch:advanced-optimizers}.
\end{important}

%------------------------------------------------------------------------------
\section{Classical Gradient Descent}
\label{sec:ch11-classical-gd}

We begin with a brief review of standard gradient descent to establish notation and highlight its limitations.

\begin{definition}{Gradient Descent Update}{ch11-gd-def}
Given a differentiable loss function $L(\vect{\theta})$, gradient descent updates parameters as:
\[
    \vect{\theta}_{t+1} = \vect{\theta}_t - \eta \grad L(\vect{\theta}_t)
\]
where $\eta > 0$ is the \textbf{learning rate}\index{learning rate|textbf}\index{step size|see{learning rate}} (or step size) and $\grad L(\vect{\theta}_t)$ is the gradient evaluated at the current parameters.
\end{definition}

\begin{intuition}
Gradient descent follows the direction of steepest descent on the loss surface. At each point, the negative gradient points toward the direction that decreases the loss most rapidly (locally). The learning rate controls how far we step in this direction---the gradient tells you \emph{where} to go, the learning rate tells you \emph{how far}. Too far and you overshoot; too short and you crawl. The optimal step depends on local curvature.
\end{intuition}

\begin{algorithmtrace}
\textbf{Example: Minimize} $f(x) = x^2$, \textbf{starting at} $x = 4$, $\eta = 0.3$

Iteration 0: $x = 4.000$, $f(x) = 16.000$, $\nabla f = 2x = 8.000$
             $x_{\text{new}} = 4.000 - 0.3 \times 8.000 = 1.600$

Iteration 1: $x = 1.600$, $f(x) = 2.560$, $\nabla f = 3.200$
             $x_{\text{new}} = 1.600 - 0.3 \times 3.200 = 0.640$

Iteration 2: $x = 0.640$, $f(x) = 0.410$, $\nabla f = 1.280$
             $x_{\text{new}} = 0.640 - 0.3 \times 1.280 = 0.256$

Iteration 3: $x = 0.256$, $f(x) = 0.066$, $\nabla f = 0.512$
             $x_{\text{new}} = 0.256 - 0.3 \times 0.512 = 0.102$

Iteration 4: $x = 0.102$, $f(x) = 0.010$, $\nabla f = 0.205$
             $x_{\text{new}} = 0.102 - 0.3 \times 0.205 = 0.041$

$\rightarrow$ Converging to $x^* = 0$ \checkmark
\end{algorithmtrace}

The optimization path is visualized in Figure~\ref{fig:gradient-descent-path}.

\begin{figure}[htbp]
    \centering
    \pdftooltip{\includegraphics[width=0.85\textwidth]{../../images/diagrams/ch11/gradient-descent-path.png}}{3D surface plot of a bowl-shaped loss landscape with a trajectory of connected dots showing the gradient descent path. The path starts at a high point on the surface and descends in steps toward the minimum at the bottom of the bowl, with step sizes decreasing as it approaches the minimum.}
    \caption{Gradient descent navigating a loss landscape. Starting from the initial point (top), the algorithm follows the negative gradient direction, taking steps toward the minimum. The path shows the characteristic convergence pattern on a convex surface.}
    \label{fig:gradient-descent-path}
\end{figure}

\subsection{Limitations of Vanilla Gradient Descent}

\begin{definition}{Condition Number}{condition-number}
For a quadratic function $f(\vect{x}) = \frac{1}{2}\vect{x}^\top \mat{H} \vect{x}$, the \textbf{condition number} is:
\[
\kappa = \frac{\lambda_{\max}(\mat{H})}{\lambda_{\min}(\mat{H})}
\]
where $\lambda_{\max}$ and $\lambda_{\min}$ are the largest and smallest eigenvalues of the Hessian $\mat{H}$.
\end{definition}

\begin{keyconcept}{Interpreting Condition Numbers}
The condition number measures how ``stretched'' the loss landscape is:
\begin{itemize}
    \item $\kappa = 1$: Perfectly spherical contours. GD converges in one step (for quadratics).
    \item $\kappa \approx 10$: Mildly elongated. GD converges reasonably fast.
    \item $\kappa \approx 100$: Highly elongated ``ravine''. GD zigzags and converges slowly.
    \item $\kappa > 1000$: Severely ill-conditioned. GD essentially stalls without preconditioning.
\end{itemize}
In practice, neural network loss landscapes often have $\kappa > 100$, which is why momentum and adaptive methods help so much.
\end{keyconcept}

\begin{example}{The Ravine Problem}{ch11-ravine}
Consider minimizing the quadratic function:
\[
    f(x, y) = \frac{1}{2}x^2 + 50y^2
\]
The Hessian has eigenvalues 1 and 100, giving condition number\index{condition number} $\kappa = 100$. Starting from $(10, 1)$ with learning rate $\eta = 0.01$:

\begin{center}
\begin{tikzpicture}
\begin{axis}[
    book figure,
    width=0.7\textwidth,
    height=0.45\textwidth,
    xlabel={$x$},
    ylabel={$y$},
    xmin=-2, xmax=12,
    ymin=-1.5, ymax=1.5,
    view={0}{90},
    colormap/viridis,
    legend pos=north east,
]
% Contour ellipses for f(x,y) = 0.5*x^2 + 50*y^2
% Level c: x^2/(2c) + y^2/(c/50) = 1, so a=sqrt(2c), b=sqrt(c/50)
\draw[gray!60, thick] (axis cs:0,0) ellipse [x radius=1.414, y radius=0.141];  % c=1
\draw[gray!60, thick] (axis cs:0,0) ellipse [x radius=3.162, y radius=0.316];  % c=5
\draw[gray!60, thick] (axis cs:0,0) ellipse [x radius=6.325, y radius=0.632];  % c=20
\draw[gray!60, thick] (axis cs:0,0) ellipse [x radius=10.0, y radius=1.0];     % c=50
\draw[gray!40] (axis cs:0,0) ellipse [x radius=14.14, y radius=1.414];         % c=100

% Gradient descent path (oscillating)
\addplot[thick, primaryred, mark=*, mark size=1pt] coordinates {
    (10,1) (9,0.0) (8.1,0.9) (7.29,-0.81) (6.561,0.729) (5.905,-0.656)
    (5.314,0.591) (4.783,-0.532) (4.305,0.478) (3.874,-0.431)
};
\addlegendentry{GD trajectory}
\node[right] at (axis cs:10,1) {\small Start};
\end{axis}
\end{tikzpicture}
\end{center}

The optimization oscillates wildly in the $y$ direction (steep) while making slow progress in the $x$ direction (shallow). This ``ravine'' or ``valley'' structure is common in neural network loss surfaces.
\end{example}

\begin{theorem}{Convergence Rate and Condition Number}{ch11-gd-convergence}
For a strongly convex quadratic $f(\vect{\theta}) = \frac{1}{2}\vect{\theta}\trans \mat{A} \vect{\theta}$ with condition number $\kappa = \lambda_{\max}(\mat{A})/\lambda_{\min}(\mat{A})$, gradient descent with optimal learning rate converges\index{convergence rate|textbf} as:
\[
    \norm{\vect{\theta}_t - \vect{\theta}^*} \leq \left(\frac{\kappa - 1}{\kappa + 1}\right)^t \norm{\vect{\theta}_0 - \vect{\theta}^*}
\]
For $\kappa = 100$, the contraction factor is approximately $0.98$, requiring about 460 iterations to reduce the error by a factor of $10^{-4}$.
\end{theorem}

This slow convergence\index{sublinear convergence} in ill-conditioned\index{ill-conditioning} problems motivates the development of accelerated methods. For rigorous convergence proofs and the mathematical foundations underlying these rates (Lipschitz continuity, strong convexity, and the descent lemma), see Chapter~\ref{ch:convergence-analysis}.

\begin{checkpoint}
What is the maximum stable learning rate for gradient descent on a function with Lipschitz constant $L$? Why?
\end{checkpoint}

%------------------------------------------------------------------------------
\section{Momentum Methods}
\label{sec:ch11-momentum}

Momentum methods address gradient descent's oscillation problem by incorporating information from past gradients, allowing the optimizer to build ``velocity'' in consistent directions.

\subsection{Classical Momentum}

\begin{definition}{SGD with Momentum}{ch11-momentum-def}
Classical momentum, introduced by Polyak~\cite{polyak1964some}, maintains a velocity vector $\vect{v}$ that accumulates gradients:\index{momentum|textbf}
\begin{align}
    \vect{v}_{t+1} &= \gamma \vect{v}_t + \grad L(\vect{\theta}_t) \\
    \vect{\theta}_{t+1} &= \vect{\theta}_t - \eta \vect{v}_{t+1}
\end{align}
where $\gamma \in [0, 1)$ is the \textbf{momentum coefficient}\index{momentum!coefficient}\index{heavy ball method|see{momentum}}\index{Polyak momentum|see{momentum}} (typically $0.9$) and $\eta$ is the learning rate.
\end{definition}

\begin{important}
\textbf{Momentum Amplifies Learning Rate.} With momentum coefficient $\gamma$, the effective learning rate becomes:
\[
\eta_{\text{effective}} = \frac{\eta}{1 - \gamma}
\]
For $\gamma = 0.9$, this is a \textbf{10x amplification}. For $\gamma = 0.99$, it's \textbf{100x}.

\textbf{Practical consequence}: When adding momentum to a working configuration, \textit{reduce your learning rate} by a factor of $(1-\gamma)$. If you were using $\eta = 0.01$ without momentum, use $\eta = 0.001$ with $\gamma = 0.9$.

\textbf{Diagnostic}: If your training diverges immediately after adding momentum, this is almost certainly why. Reduce the learning rate by 5--10x.
\end{important}

\begin{important}
\textbf{Why 0.9 is the Default:} The momentum value 0.9 appears everywhere---PyTorch, TensorFlow, and virtually every deep learning tutorial. Empirical practice has converged on 0.9 as a robust default: it provides enough smoothing to escape local noise (effective memory of $\sim$10 steps) while remaining responsive to gradient direction changes. At $\gamma=0.99$, the 100-step memory makes the optimizer too sluggish; at $\gamma=0.8$, the 5-step memory provides insufficient smoothing. Unless you have strong evidence otherwise, use $\gamma = 0.9$---decades of empirical work have validated this default.
\end{important}

\begin{intuition}
Imagine a ball rolling down a hilly terrain. Without momentum, the ball would stop immediately when encountering a flat region or upward slope. With momentum, the ball accumulates velocity as it rolls, allowing it to push through small bumps and oscillations. The momentum coefficient $\gamma$ determines how much ``inertia'' the ball carries---higher values mean longer memory of past motion.
\end{intuition}

The effective update can be written as an exponentially weighted sum of past gradients:
\[
    \vect{v}_t = \sum_{i=0}^{t} \gamma^{t-i} \grad L(\vect{\theta}_i)
\]

This reveals that momentum gives more weight to recent gradients while gradually forgetting older ones with decay factor $\gamma^{t-i}$.

\begin{theorem}{Momentum Acceleration}{ch11-momentum-accel}
For convex quadratic objectives with condition number $\kappa$, gradient descent with optimal momentum achieves\index{convergence rate!linear}\index{linear convergence}:
\[
    \norm{\vect{\theta}_t - \vect{\theta}^*} \leq \left(\frac{\sqrt{\kappa} - 1}{\sqrt{\kappa} + 1}\right)^t \norm{\vect{\theta}_0 - \vect{\theta}^*}
\]
This improves upon standard gradient descent by effectively replacing $\kappa$ with $\sqrt{\kappa}$. For $\kappa = 100$, the contraction factor improves from $0.98$ to approximately $0.82$.
\end{theorem}

\begin{checkpoint}
How does momentum help in narrow valleys? Draw the trajectory with and without momentum.
\end{checkpoint}

While momentum accelerates convergence, it has one limitation: it can overshoot near the optimum, oscillating back and forth. Nesterov's key innovation addresses this by ``looking ahead'' before committing to a step.

\subsection{Nesterov Accelerated Gradient}

\begin{definition}{Nesterov Momentum}{ch11-nesterov-def}
\textbf{Nesterov Accelerated Gradient}\index{Nesterov momentum|textbf}\index{accelerated gradient descent|see{Nesterov momentum}}\index{NAG|see{Nesterov momentum}} (NAG) evaluates the gradient at a ``look-ahead'' position:
\begin{align}
    \vect{v}_{t+1} &= \gamma \vect{v}_t + \grad L(\vect{\theta}_t - \eta \gamma \vect{v}_t) \\
    \vect{\theta}_{t+1} &= \vect{\theta}_t - \eta \vect{v}_{t+1}
\end{align}
The key difference from classical momentum is evaluating the gradient at $\vect{\theta}_t - \eta \gamma \vect{v}_t$ rather than $\vect{\theta}_t$.

\textbf{Note:} This formulation is the momentum-based ``look-ahead'' variant commonly used in deep learning (Sutskever et al., 2013), which differs from Nesterov's original accelerated gradient method (Nesterov, 1983). The original method was derived for convex optimization with specific step size schedules to achieve optimal $O(1/k^2)$ convergence rates. The momentum-based variant presented here adapts the look-ahead idea for practical deep learning, where it provides similar benefits of anticipatory correction without requiring the theoretical step size schedules.
\end{definition}

\begin{intuition}
Classical momentum computes the gradient at the current position, then applies the momentum step. Nesterov reverses this: first take a provisional step in the momentum direction, then compute the gradient there. This ``look-ahead'' provides early warning---if momentum is about to overshoot, the gradient at the look-ahead position will point backward, applying a correction before the mistake happens. Intuitively, Nesterov momentum is like a skier who looks ahead to see the terrain before turning---if there's a cliff coming, they start turning early instead of waiting until they're at the edge.
\end{intuition}

\begin{keyconcept}
\textbf{Classical Momentum vs. Nesterov Momentum---A Direct Comparison:}

\begin{center}
\begin{tabular}{lll}
\toprule
\textbf{Aspect} & \textbf{Classical (Polyak)} & \textbf{Nesterov (NAG)} \\
\midrule
Gradient evaluation & At current position $\vect{\theta}_t$ & At look-ahead $\vect{\theta}_t - \eta\gamma\vect{v}_t$ \\
Correction timing & After overshooting & Before overshooting \\
Convergence rate & $O(1/\sqrt{\kappa})$ & $O(1/\sqrt{\kappa})$ (with better constants) \\
Near optima & May overshoot, then correct & Anticipates and slows down \\
\bottomrule
\end{tabular}
\end{center}

\textbf{When does Nesterov help most?}
\begin{itemize}
    \item Near the optimum, where overshooting is costly
    \item On ill-conditioned problems with narrow valleys
    \item When momentum is high ($\gamma \geq 0.9$)
\end{itemize}

\textbf{Practical reality:} For well-tuned hyperparameters, the difference between classical and Nesterov momentum is often small---perhaps 5--10\% fewer iterations to converge. Both dramatically outperform vanilla gradient descent. If you are using an adaptive optimizer like Adam (which has its own momentum via $\beta_1$), the classical vs. Nesterov distinction becomes less important. Focus on Nesterov when using SGD with momentum on convex or mildly non-convex problems.
\end{keyconcept}

\begin{checkpoint}
What's the difference between Polyak momentum and Nesterov momentum? Which looks ahead?
\end{checkpoint}

\begin{important}
\textbf{Ready for Adaptive Methods?} This chapter covered gradient descent and momentum methods---the foundations. For adaptive learning rate methods (AdaGrad, RMSprop, Adam, AdamW) that automatically tune per-parameter learning rates, continue to Chapter~\ref{ch:advanced-optimizers}.
\end{important}

%------------------------------------------------------------------------------
\section{Worked Example: Optimizer Comparison on the Rosenbrock Function}
\label{sec:ch11-optimizer-showdown}

\begin{important}
\textbf{Preview of Coming Concepts.} This worked example compares SGD+Momentum with RMSprop and Adam optimizers. RMSprop and Adam are fully introduced in Chapter~\ref{ch:advanced-optimizers}. For now, focus on comparing the trajectory behaviors and convergence characteristics; the algorithm details will be explained in the next chapter.
\end{important}

To build practical intuition about optimizer behavior, we now conduct a systematic comparison on a challenging test function.

\subsection{The Rosenbrock Function}

\begin{definition}{Rosenbrock Function}{ch11-rosenbrock-def}
The \textbf{Rosenbrock function}\index{Rosenbrock function|textbf} is a classic optimization test case:
\[
    f(x, y) = (1 - x)^2 + 100(y - x^2)^2
\]
The global minimum is at $(x^*, y^*) = (1, 1)$ where $f(1, 1) = 0$.
\end{definition}

The Rosenbrock function presents a banana-shaped valley that is notoriously difficult to optimize. The valley floor has a gentle slope toward the minimum, but the valley walls are steep. This geometry creates challenges similar to those encountered in neural network optimization:

\begin{itemize}
    \item \textbf{Ill-conditioning}: The Hessian condition number varies dramatically across the domain
    \item \textbf{Curved valley}: The optimal direction changes along the trajectory
    \item \textbf{Deceptive gradients}: The gradient often points across the valley rather than along it
\end{itemize}

The gradient is:
\begin{align}
    \pd{f}{x} &= -2(1-x) - 400x(y - x^2) \\
    \pd{f}{y} &= 200(y - x^2)
\end{align}

\subsection{Experimental Setup}

We compare three optimizers starting from $\vect{\theta}_0 = (-1, 1)$ with the following hyperparameters:

\begin{table}[htbp]
\centering
\caption{Hyperparameters for optimizer comparison experiment. (RMSprop and Adam are introduced in Chapter~\ref{ch:advanced-optimizers}.)}
\label{tab:ch11-hyperparams}
\begin{tabular}{@{}llp{6cm}@{}}
\toprule
\textbf{Optimizer} & \textbf{Hyperparameters} & \textbf{Notes} \\
\midrule
SGD + Momentum & $\eta = 0.0005$, $\gamma = 0.9$ & Learning rate tuned to prevent divergence \\
RMSprop & $\eta = 0.01$, $\beta = 0.9$, $\epsilon = 10^{-8}$ & Standard defaults \\
Adam & $\eta = 0.01$, $\beta_1 = 0.9$, $\beta_2 = 0.999$, $\epsilon = 10^{-8}$ & Standard defaults (slightly higher $\eta$) \\
\bottomrule
\end{tabular}
\end{table}

\begin{important}
These hyperparameters were selected to give each optimizer a fair chance. SGD requires a much smaller learning rate to remain stable, while adaptive methods tolerate larger rates. This difference in tuning sensitivity is itself an important practical consideration.
\end{important}

The resulting optimization trajectories are shown in Figure~\ref{fig:ch11-rosenbrock-trajectories}, and the corresponding convergence curves appear in Figure~\ref{fig:ch11-convergence-curves}.

\subsection{Optimization Trajectories}

\begin{figure}[htbp]
\centering
\begin{tikzpicture}
\begin{axis}[
    book figure,
    width=0.95\textwidth,
    height=0.65\textwidth,
    xlabel={$x$},
    ylabel={$y$},
    xmin=-1.5, xmax=1.5,
    ymin=-0.5, ymax=1.5,
    view={0}{90},
    colorbar,
    colorbar style={ylabel={$\ln(f(x,y) + 1)$}},
    point meta min=0,
    point meta max=3,
    legend pos=outer north east,
]

% Surface plot of Rosenbrock function (no gnuplot required)
\addplot3[
    surf,
    shader=interp,
    domain=-1.5:1.5,
    domain y=-0.5:1.5,
    samples=40,
    opacity=0.7,
] {ln((1-x)^2 + 100*(y-x^2)^2 + 1)};
% Valley curve y = x^2 shown as guide
\addplot[black!40, dashed, thick, domain=-1.2:1.2, samples=50] ({x}, {x^2});

% Mark the minimum
\addplot[only marks, mark=star, mark size=4pt, primarygreen!80!black] coordinates {(1,1)};
\node[above right, primarygreen!80!black] at (axis cs:1,1) {\small Minimum $(1,1)$};

% SGD + Momentum trajectory (blue) - oscillates, slow progress
\addplot[thick, primaryblue, mark=*, mark size=0.8pt, mark repeat=50] coordinates {
    (-1.0, 1.0) (-0.92, 0.85) (-0.84, 0.72) (-0.77, 0.61) (-0.71, 0.52)
    (-0.65, 0.44) (-0.60, 0.38) (-0.55, 0.33) (-0.51, 0.29) (-0.47, 0.25)
    (-0.43, 0.22) (-0.39, 0.19) (-0.36, 0.17) (-0.33, 0.15) (-0.30, 0.13)
    (-0.27, 0.11) (-0.24, 0.10) (-0.21, 0.08) (-0.18, 0.07) (-0.15, 0.06)
    (-0.12, 0.05) (-0.09, 0.04) (-0.06, 0.03) (-0.03, 0.02) (0.00, 0.02)
    (0.03, 0.01) (0.06, 0.01) (0.09, 0.01) (0.12, 0.02) (0.15, 0.03)
    (0.18, 0.04) (0.21, 0.05) (0.24, 0.06) (0.27, 0.08) (0.30, 0.10)
    (0.33, 0.12) (0.36, 0.14) (0.39, 0.16) (0.42, 0.18) (0.45, 0.21)
    (0.48, 0.24) (0.51, 0.27) (0.54, 0.30) (0.57, 0.33) (0.60, 0.37)
    (0.63, 0.40) (0.66, 0.44) (0.69, 0.48) (0.72, 0.52) (0.75, 0.57)
};

% RMSprop trajectory (red) - faster, follows valley
\addplot[thick, primaryred, mark=square*, mark size=0.8pt, mark repeat=20] coordinates {
    (-1.0, 1.0) (-0.85, 0.78) (-0.71, 0.58) (-0.58, 0.42) (-0.46, 0.30)
    (-0.35, 0.21) (-0.24, 0.14) (-0.14, 0.09) (-0.05, 0.05) (0.04, 0.03)
    (0.13, 0.03) (0.22, 0.05) (0.30, 0.10) (0.38, 0.15) (0.45, 0.21)
    (0.52, 0.28) (0.58, 0.35) (0.64, 0.42) (0.70, 0.49) (0.75, 0.57)
    (0.80, 0.64) (0.84, 0.71) (0.88, 0.78) (0.91, 0.84) (0.94, 0.89)
    (0.96, 0.93) (0.98, 0.96) (0.99, 0.98) (1.00, 0.99) (1.00, 1.00)
};

% Adam trajectory (green) - fastest, smoothest
\addplot[thick, primarygreen, mark=triangle*, mark size=0.8pt, mark repeat=15] coordinates {
    (-1.0, 1.0) (-0.82, 0.72) (-0.65, 0.50) (-0.49, 0.34) (-0.34, 0.22)
    (-0.20, 0.14) (-0.07, 0.08) (0.05, 0.05) (0.16, 0.04) (0.27, 0.08)
    (0.37, 0.14) (0.46, 0.22) (0.55, 0.31) (0.63, 0.40) (0.70, 0.50)
    (0.77, 0.60) (0.83, 0.69) (0.88, 0.78) (0.92, 0.85) (0.95, 0.91)
    (0.97, 0.95) (0.99, 0.98) (1.00, 0.99) (1.00, 1.00)
};

% Start point
\addplot[only marks, mark=o, mark size=4pt, black, thick] coordinates {(-1,1)};
\node[above left] at (axis cs:-1,1) {\small Start $(-1,1)$};

% Legend
\addlegendentry{SGD+Momentum}
\addlegendentry{RMSprop}
\addlegendentry{Adam}

\end{axis}
\end{tikzpicture}
\caption{Optimization trajectories on the Rosenbrock function. Contour lines show the loss surface with the banana-shaped valley. Adam (green) reaches the minimum fastest, followed by RMSprop (red). SGD with momentum (blue) makes slower progress and requires significantly more iterations. \textbf{Note:} This is a schematic illustration; actual trajectories depend on specific hyperparameter choices and random initialization.}
\label{fig:ch11-rosenbrock-trajectories}
\end{figure}

\subsection{Convergence Analysis}

\begin{figure}[htbp]
\centering
\begin{tikzpicture}
\begin{axis}[
    book figure,
    width=0.85\textwidth,
    height=0.45\textwidth,
    xlabel={Iterations},
    ylabel={$\log_{10}(f(\vect{\theta}_t) + 10^{-10})$},
    xmin=0, xmax=5000,
    ymin=-8, ymax=4,
    legend pos=north east,
    legend cell align={left},
]

% SGD + Momentum (slow decline, plateau)
\addplot[thick, primaryblue, domain=0:5000, samples=100]
    {3.5 * exp(-0.0008*x) + 0.5*exp(-0.0002*x) - 2};
\addlegendentry{SGD+Momentum}

% RMSprop (faster, reaches lower)
\addplot[thick, primaryred, domain=0:5000, samples=100]
    {3.5 * exp(-0.003*x) - 5};
\addlegendentry{RMSprop}

% Adam (fastest convergence)
\addplot[thick, primarygreen, domain=0:5000, samples=100]
    {3.5 * exp(-0.005*x) - 7};
\addlegendentry{Adam}

% Convergence threshold epsilon = 0.01
\addplot[dashed, black, domain=0:5000] {-2};
\node[right] at (axis cs:5000,-2) {\small $\varepsilon = 0.01$};

\end{axis}
\end{tikzpicture}
\caption{Convergence curves showing loss versus iteration count. The dashed line marks $\varepsilon = 0.01$ convergence threshold. Adam reaches this threshold fastest, followed by RMSprop, with SGD+Momentum requiring substantially more iterations. \textbf{Note:} Curves are schematic to illustrate typical relative behavior; actual convergence depends on hyperparameters.}
\label{fig:ch11-convergence-curves}
\end{figure}

\subsection{Quantitative Comparison}

\begin{table}[htbp]
\centering
\small
\caption{Optimizer performance comparison on the Rosenbrock function. Iteration counts are representative of typical behavior with the hyperparameters in Table~\ref{tab:ch11-hyperparams}; actual values vary with initialization and tuning.}
\label{tab:ch11-optimizer-comparison}
\begin{tabular}{@{}lccc>{\raggedright\arraybackslash}p{5.5cm}@{}}
\toprule
\textbf{Optimizer} & \textbf{Iters to $\varepsilon$\textsuperscript{*}} & \textbf{Tuning} & \textbf{Memory} & \textbf{Best For} \\
\midrule
SGD+Momentum & $\sim$5000 & High & $O(d)$ & Convex, well-conditioned problems; CNNs with careful tuning \\
\addlinespace
RMSprop & $\sim$800 & Medium & $O(d)$ & Non-stationary objectives; RNNs; online learning \\
\addlinespace
Adam & $\sim$500 & Low & $O(2d)$ & General deep learning; quick prototyping; Transformers \\
\bottomrule
\multicolumn{5}{@{}l}{\textsuperscript{*}\footnotesize Values shown are representative of typical behavior, not from a specific experiment.}
\end{tabular}
\end{table}

\begin{keyconcept}
\textbf{Key Patterns Illustrated by This Comparison:}

\begin{enumerate}
    \item \textbf{Adam typically converges fastest} with minimal tuning---the default hyperparameters tend to work well out of the box. This characteristic makes it well-suited for rapid prototyping and new projects.

    \item \textbf{SGD+Momentum requires careful learning rate selection}. Too large causes divergence; too small causes unacceptably slow convergence. However, with proper tuning, it can match or exceed Adam's final solution quality.

    \item \textbf{RMSprop offers a middle ground}---generally faster than SGD, more interpretable than Adam, and particularly effective for recurrent architectures where gradient magnitudes vary significantly.

    \item \textbf{Trajectory shapes characteristically differ}: SGD tends to follow a more oscillatory path, while adaptive methods typically navigate curved valleys more smoothly by automatically adjusting step sizes in different directions.
\end{enumerate}

\textbf{Note:} These patterns reflect typical optimizer behavior observed across many problems; actual performance varies with hyperparameters and problem structure.
\end{keyconcept}

\subsection{Implementation}

\begin{pythoncode}
import numpy as np

def rosenbrock(theta):
    """Rosenbrock function: f(x,y) = (1-x)^2 + 100(y-x^2)^2"""
    x, y = theta
    return (1 - x)**2 + 100 * (y - x**2)**2

def rosenbrock_grad(theta):
    """Gradient of Rosenbrock function."""
    x, y = theta
    dx = -2 * (1 - x) - 400 * x * (y - x**2)
    dy = 200 * (y - x**2)
    return np.array([dx, dy])

def sgd_momentum(grad_fn, theta_init, lr=0.0005, momentum=0.9, max_iters=5000):
    """SGD with momentum optimizer."""
    theta = theta_init.copy()
    velocity = np.zeros_like(theta)
    history = [theta.copy()]

    for _ in range(max_iters):
        grad = grad_fn(theta)
        velocity = momentum * velocity + grad
        theta = theta - lr * velocity
        history.append(theta.copy())

    return theta, np.array(history)

def rmsprop(grad_fn, theta_init, lr=0.01, beta=0.9, eps=1e-8, max_iters=5000):
    """RMSprop optimizer."""
    theta = theta_init.copy()
    v = np.zeros_like(theta)
    history = [theta.copy()]

    for _ in range(max_iters):
        grad = grad_fn(theta)
        v = beta * v + (1 - beta) * grad**2
        theta = theta - lr * grad / (np.sqrt(v) + eps)
        history.append(theta.copy())

    return theta, np.array(history)

def adam(grad_fn, theta_init, lr=0.01, beta1=0.9, beta2=0.999,
         eps=1e-8, max_iters=5000):
    """Adam optimizer."""
    theta = theta_init.copy()
    m, v = np.zeros_like(theta), np.zeros_like(theta)
    history = [theta.copy()]

    for t in range(1, max_iters + 1):
        grad = grad_fn(theta)
        m = beta1 * m + (1 - beta1) * grad
        v = beta2 * v + (1 - beta2) * grad**2
        m_hat = m / (1 - beta1**t)
        v_hat = v / (1 - beta2**t)
        theta = theta - lr * m_hat / (np.sqrt(v_hat) + eps)
        history.append(theta.copy())

    return theta, np.array(history)

# Run comparison
theta_init = np.array([-1.0, 1.0])
theta_sgd, hist_sgd = sgd_momentum(rosenbrock_grad, theta_init)
theta_rms, hist_rms = rmsprop(rosenbrock_grad, theta_init)
theta_adam, hist_adam = adam(rosenbrock_grad, theta_init)

# Count iterations to epsilon-convergence
def iters_to_converge(history, loss_fn, eps=0.01):
    for i, theta in enumerate(history):
        if loss_fn(theta) < eps:
            return i
    return len(history)

print(f"SGD+Momentum: {iters_to_converge(hist_sgd, rosenbrock)} iters")
print(f"RMSprop: {iters_to_converge(hist_rms, rosenbrock)} iters")
print(f"Adam: {iters_to_converge(hist_adam, rosenbrock)} iters")
\end{pythoncode}

%------------------------------------------------------------------------------
\section{Diagnosing Momentum and Gradient Descent Issues}
\label{sec:ch11-diagnosing}

This section provides systematic guidance for diagnosing and fixing momentum-related training problems.

\subsection{Oscillation Patterns and What They Mean}

\begin{important}
\textbf{Diagnosing Oscillation Patterns:}

\textbf{Oscillations with increasing amplitude:} Momentum is amplifying instability from too-large learning rate. With $\gamma = 0.9$, the effective step size is 10x the base rate. \emph{Diagnosis:} Temporarily set $\gamma = 0$; if oscillations stop, momentum is the culprit. \emph{Fix:} Reduce learning rate by factor of $(1 - \gamma)$ (i.e., 10x for $\gamma = 0.9$), or reduce momentum to $\gamma = 0.5$.

\textbf{Oscillations with constant amplitude:} Momentum is at the stability boundary---not diverging, but not converging. \emph{Diagnosis:} Check if reducing $\eta$ by 20\% causes oscillations to shrink. \emph{Fix:} Reduce learning rate by 30--50\%, add learning rate decay, or switch to Adam.

\textbf{Example:} With $\eta = 0.01$ and $\gamma = 0.9$, the effective learning rate is $\eta/(1-\gamma) = 0.1$---a 10$\times$ amplification. If training was stable at $\eta = 0.1$ without momentum, adding $\gamma = 0.9$ requires reducing to $\eta = 0.01$.
\end{important}

\begin{keyconcept}
\textbf{Visual Signatures of Momentum Problems:}

\textbf{Healthy momentum trajectory:}
\begin{itemize}
    \item Smooth, curved path toward minimum
    \item May overshoot slightly, then correct
    \item Converges faster than no-momentum baseline
\end{itemize}

\textbf{Too much momentum ($\gamma$ too high):}
\begin{itemize}
    \item Wild oscillations that grow over time
    \item Trajectory ``spirals outward'' around the minimum
    \item Loss curve shows repeated spikes
\end{itemize}

\textbf{Momentum + high learning rate:}
\begin{itemize}
    \item Rapid initial progress followed by explosion
    \item Parameters may hit NaN after appearing to converge
    \item Velocity accumulates to dangerous levels
\end{itemize}

\textbf{Momentum on noisy gradients (when it hurts):}
\begin{itemize}
    \item Loss oscillates more with momentum than without
    \item Velocity accumulates noise instead of signal
    \item Common in RNNs with long sequences or very small batch sizes
\end{itemize}
\end{keyconcept}

\subsection{When Momentum Hurts: Noisy Gradient Settings}

\begin{commonmistakes}
\begin{itemize}
    \item \textbf{Mistake:} Always using momentum because ``it accelerates convergence''

    \textbf{Why it's wrong:} Momentum accelerates convergence when gradients are \textit{consistent}---pointing in similar directions across steps. With highly noisy gradients (small batch, high variance), momentum accumulates noise as much as signal.

    \textbf{Correct understanding:} For very small batches (1--8 samples) or tasks with inherently noisy gradients (e.g., reinforcement learning, GANs), momentum can hurt. Try reducing $\gamma$ to 0.5 or 0, or use adaptive methods (Adam) that handle noise better.
\end{itemize}
\end{commonmistakes}

\begin{principle}
\textbf{When to Reduce or Remove Momentum:}

\begin{enumerate}
    \item \textbf{Very small batch sizes} ($B < 16$): Gradient noise dominates; momentum amplifies noise
    \item \textbf{RNNs with long sequences}: Gradients vary dramatically across time; momentum on the wrong signal
    \item \textbf{GANs}: Non-stationary optimization; momentum can destabilize the adversarial game
    \item \textbf{Early training with random initialization}: Gradients point in random directions; momentum averages noise
    \item \textbf{Fine-tuning pretrained models}: Gradients should be small; momentum can overshoot the nearby optimum
\end{enumerate}

In these cases, use $\gamma = 0$ to $0.5$, or switch to Adam (which has implicit dampening via adaptive rates).
\end{principle}

\subsection{Diagnosing Overshooting}

\begin{important}
\textbf{Symptom: Loss decreases rapidly, then suddenly increases and stays high}

\textbf{Likely cause:} Momentum built up velocity toward the minimum, then overshot and landed in a bad region.

\textbf{Why it happens:} Near a minimum, gradients become small. But momentum $\vect{v}_t = \gamma \vect{v}_{t-1} + \vect{g}_t$ still carries large velocity from earlier steps. The small corrective gradient cannot stop the large momentum, causing overshoot.

\textbf{Diagnosis:}
\begin{enumerate}
    \item Track $\norm{\vect{v}_t}$ over time---if it peaks just before the loss spike, overshoot is confirmed
    \item Check if $\norm{\vect{v}_t} \gg \norm{\vect{g}_t}$ near the spike (velocity dominates gradient)
\end{enumerate}

\textbf{Fixes:}
\begin{enumerate}
    \item Use learning rate decay to reduce step size as you approach the minimum
    \item Use Nesterov momentum (look-ahead gradient provides earlier warning)
    \item Reduce momentum coefficient to $\gamma = 0.8$ or lower
    \item Add gradient clipping to bound $\norm{\vect{g}_t}$ (limits velocity accumulation)
\end{enumerate}
\end{important}

\begin{warstory}
\textbf{The Momentum That Killed a GAN}

A team training a StyleGAN variant noticed the model generated beautiful images for 100k iterations, then suddenly collapsed to producing noise. Loss plots showed the generator loss creeping up before exploding.

The diagnosis: momentum of $\gamma = 0.9$ was accumulating velocity in the generator's winning direction against the discriminator. When the discriminator finally adapted, the generator had so much momentum that it couldn't change direction fast enough, overshooting into a degenerate region.

The fix was two-fold: (1) reduce generator momentum to $\gamma = 0.5$ while keeping discriminator at $\gamma = 0.9$, and (2) add exponential moving average of generator weights for evaluation. The asymmetric momentum prevented the generator from ``running away'' while still allowing the discriminator to learn quickly.

\textbf{Lesson:} In adversarial training, momentum interacts with the game dynamics. One player's acceleration becomes another player's instability. Tune momentum per-network, not globally.
\end{warstory}

\subsection{Learning Rate and Momentum Interaction}

\begin{keyconcept}
\textbf{The Stability Triangle:}

For SGD with momentum, stability requires:
\[
    \eta < \frac{2(1 + \gamma)}{L}
\]
where $L$ is the smoothness constant. Key implications:

\begin{itemize}
    \item Higher momentum ($\gamma \uparrow$) allows higher learning rate for stability
    \item But effective step size is $\eta_{\text{eff}} \approx \eta / (1 - \gamma)$
    \item Net effect: $\eta_{\text{eff}} < \frac{2(1 + \gamma)}{L(1 - \gamma)}$
\end{itemize}

\textbf{Practical guidance:}
\begin{itemize}
    \item If stable with $\eta = 0.1$ and $\gamma = 0$, then with $\gamma = 0.9$ use $\eta \leq 0.02$
    \item When increasing $\gamma$, decrease $\eta$ roughly proportionally
    \item Test stability: if loss oscillates, you are above the threshold
\end{itemize}
\end{keyconcept}

\begin{important}
\textbf{Quick Diagnostic Checklist for Momentum:}

\begin{enumerate}
    \item \textbf{Loss oscillates?}
    \begin{itemize}
        \item Set $\gamma = 0$ temporarily. If oscillations stop $\rightarrow$ reduce $\eta$ by $2\times$--$5\times$
        \item If oscillations persist without momentum $\rightarrow$ problem is $\eta$ alone, not momentum
    \end{itemize}

    \item \textbf{Training slower with momentum than without?}
    \begin{itemize}
        \item Check batch size: if $B < 16$, try $\gamma = 0$ or $\gamma = 0.5$
        \item Check gradient variance: if $\Var[\vect{g}] > \E[\vect{g}]^2$, momentum amplifies noise
    \end{itemize}

    \item \textbf{Loss explodes after initial progress?}
    \begin{itemize}
        \item Classic overshoot: add learning rate decay
        \item Track velocity norm: if $\norm{\vect{v}_t}$ grows unboundedly, reduce $\gamma$
    \end{itemize}

    \item \textbf{Convergence is bumpy near the end?}
    \begin{itemize}
        \item Momentum is too aggressive for fine-tuning phase
        \item Use cosine annealing or step decay to reduce effective momentum late in training
    \end{itemize}
\end{enumerate}
\end{important}

%------------------------------------------------------------------------------
\section{Exercises}
\label{sec:ch11-exercises}

\begin{exercise}{Momentum Dynamics}{ch11-ex-momentum}
Consider minimizing $f(x) = \frac{1}{2}x^2$ with SGD+Momentum starting from $x_0 = 10$.
\begin{enumerate}[label=(\alph*)]
    \item With $\eta = 0.1$ and $\gamma = 0$, derive the sequence $x_t$ and show it converges geometrically.
    \item With $\eta = 0.1$ and $\gamma = 0.9$, derive the update equations. What is the effective learning rate after the velocity builds up?
    \item At what value of $\gamma$ does the system become unstable (updates grow unboundedly)?
\end{enumerate}
\end{exercise}

\begin{exercise}{Adam Bias Correction}{ch11-ex-adam-bias}
\begin{enumerate}[label=(\alph*)]
    \item For constant gradients $\vect{g}_t = \vect{g}$ for all $t$, derive the exact value of $\vect{m}_t$ (without bias correction).
    \item Show that $\E[\vect{m}_t] = (1 - \beta_1^t) \vect{g}$, confirming the bias toward zero.
    \item Plot the bias correction factor $1/(1-\beta_1^t)$ for $t = 1, \ldots, 100$ with $\beta_1 = 0.9$. At what iteration is the correction less than 1\%?
\end{enumerate}
\end{exercise}

\begin{exercise}{Optimizer Comparison}{ch11-ex-comparison}
Implement SGD, SGD+Momentum, RMSprop, and Adam from scratch. Test them on:
\begin{enumerate}[label=(\alph*)]
    \item The Rosenbrock function starting from $(-1, 1)$
    \item The Beale function: $f(x,y) = (1.5 - x + xy)^2 + (2.25 - x + xy^2)^2 + (2.625 - x + xy^3)^2$
    \item A simple quadratic with condition number 100: $f(x,y) = x^2 + 100y^2$
\end{enumerate}
Compare iterations to convergence and discuss which optimizer performs best on each problem and why.
\end{exercise}

\begin{exercise}{Learning Rate Sensitivity}{ch11-ex-lr-sensitivity}
Using the Rosenbrock function:
\begin{enumerate}[label=(\alph*)]
    \item For SGD+Momentum ($\gamma = 0.9$), find the largest stable learning rate $\eta_{\max}$.
    \item For Adam (default $\beta$ values), find the largest stable learning rate.
    \item What is the ratio $\eta_{\max}^{\text{Adam}} / \eta_{\max}^{\text{SGD}}$? What does this tell you about tuning sensitivity?
\end{enumerate}
\end{exercise}

%------------------------------------------------------------------------------
\section{Summary}
\label{sec:ch11-summary}

This chapter covered the foundational gradient descent methods and momentum techniques:

\begin{itemize}
    \item \textbf{Vanilla gradient descent} suffers from slow convergence on ill-conditioned problems, with convergence rate depending on the condition number $\kappa$.

    \item \textbf{Momentum methods}\index{momentum!methods} accumulate gradient information over time, accelerating convergence\index{convergence} in consistent directions while dampening oscillations. Classical momentum improves the condition number dependence from $\kappa$ to $\sqrt{\kappa}$.

    \item \textbf{Nesterov momentum} provides additional improvement through look-ahead gradient evaluation, enabling earlier correction when about to overshoot.

    \item \textbf{The Rosenbrock function comparison} illustrated how momentum helps navigate ill-conditioned landscapes, though SGD+Momentum requires careful learning rate tuning compared to adaptive methods.

    \item \textbf{Diagnosing momentum issues} requires understanding the interaction between learning rate and momentum coefficient---the effective learning rate is $\eta/(1-\gamma)$, which is 10x the base rate for $\gamma=0.9$.

    \item \textbf{Learning rate remains the most critical hyperparameter}---proper tuning of the learning rate is essential for stable and efficient convergence.
\end{itemize}

For adaptive learning rate methods (AdaGrad, RMSprop, Adam, AdamW), learning rate schedules, and practical optimizer selection guidelines, continue to Chapter~\ref{ch:advanced-optimizers}.

\section*{Interview Questions}
\begin{interviewquestions}

\textbf{Question 1: How do you choose a learning rate?}

\textit{Expected follow-ups:}
\begin{itemize}
    \item What happens if the learning rate is too high? Too low?
    \item How do learning rate schedules help?
    \item Should you use the same learning rate for Adam and SGD?
\end{itemize}

\textit{Red flags to avoid:}
\begin{itemize}
    \item Giving a single number without context (``I always use 0.001'')
    \item Not mentioning that different optimizers need different learning rates
    \item Forgetting about learning rate warmup for large models
\end{itemize}

\textit{Strong answer key points:}
\begin{itemize}
    \item \textbf{Theory}: Maximum stable LR is $\eta \leq 1/L$ where $L$ is the smoothness constant (Lipschitz constant of gradient)
    \item \textbf{Practice}: Use learning rate finder---sweep LR exponentially, plot loss, choose LR where loss decreases most steeply (just before it explodes)
    \item \textbf{Defaults by optimizer}: SGD: 0.01-0.1, Adam: 0.001, AdamW: 0.001. These differ by 10-100x!
    \item \textbf{Schedules}: Start high for exploration, decay for fine convergence. Cosine annealing and step decay are most common
    \item \textbf{Warmup}: Essential for Transformers and large batch training---start small, ramp up over first 1-10\% of training
\end{itemize}

\vspace{0.5em}
\textbf{Question 2: Explain momentum and why it accelerates convergence.}

\textit{Expected follow-ups:}
\begin{itemize}
    \item What's the difference between Polyak momentum and Nesterov momentum?
    \item How does momentum help with ill-conditioned problems?
    \item What is the typical value for the momentum coefficient?
\end{itemize}

\textit{Red flags to avoid:}
\begin{itemize}
    \item Saying momentum just ``averages gradients''---it accumulates velocity
    \item Not knowing that momentum improves condition number dependence from $\kappa$ to $\sqrt{\kappa}$
    \item Confusing momentum's smoothing effect with reducing gradient variance
\end{itemize}

\textit{Strong answer key points:}
\begin{itemize}
    \item \textbf{Intuition}: Ball rolling down a hill accumulates velocity, pushing through small bumps and oscillations
    \item \textbf{Update}: $v_{t+1} = \gamma v_t + \nabla L(\theta_t)$, then $\theta_{t+1} = \theta_t - \eta v_{t+1}$
    \item \textbf{Acceleration}: Improves convergence from $O((1-1/\kappa)^t)$ to $O((1-1/\sqrt{\kappa})^t)$---huge for ill-conditioned problems
    \item \textbf{Nesterov}: Evaluates gradient at ``look-ahead'' position $\theta - \eta\gamma v$, providing early correction if about to overshoot
    \item \textbf{Typical value}: $\gamma = 0.9$ works well in most cases; 0.99 for very noisy gradients
\end{itemize}

\vspace{0.5em}
\textbf{Question 3: What causes oscillation in gradient descent and how do you fix it?}

\textit{Expected follow-ups:}
\begin{itemize}
    \item How does the condition number relate to oscillation?
    \item Which optimizers are less prone to oscillation?
    \item How can you diagnose oscillation during training?
\end{itemize}

\textit{Red flags to avoid:}
\begin{itemize}
    \item Only suggesting ``reduce learning rate'' without explaining why
    \item Not mentioning ill-conditioning as the root cause
    \item Forgetting that adaptive methods naturally handle this
\end{itemize}

\textit{Strong answer key points:}
\begin{itemize}
    \item \textbf{Root cause}: High condition number $\kappa = L/\mu$---the loss landscape is ``elongated'' like a narrow valley
    \item \textbf{Mechanism}: Large gradients in steep directions cause overshooting; small gradients in flat directions cause slow progress
    \item \textbf{Solutions}: (1) Reduce learning rate (reduces oscillation but also progress), (2) Add momentum (dampens oscillations while maintaining velocity), (3) Use adaptive optimizers (Adam, RMSprop) that normalize per-parameter
    \item \textbf{Diagnosis}: Plot loss over iterations---oscillation appears as sawtooth pattern that doesn't decrease
    \item \textbf{Preconditioning}: Batch normalization effectively reduces condition number, enabling larger learning rates
\end{itemize}

\vspace{0.5em}
\textbf{Question 4: Compare gradient descent to Newton's method. When would you use each?}

\textit{Expected follow-ups:}
\begin{itemize}
    \item What is the computational complexity of Newton's method?
    \item Are there approximations to Newton's method used in deep learning?
    \item Why isn't Newton's method used for training neural networks?
\end{itemize}

\textit{Red flags to avoid:}
\begin{itemize}
    \item Not knowing that Newton's method uses second-order (Hessian) information
    \item Saying Newton's method is ``always better'' without mentioning computational cost
    \item Not connecting adaptive methods to approximate second-order information
\end{itemize}

\textit{Strong answer key points:}
\begin{itemize}
    \item \textbf{Gradient descent}: First-order, $O(d)$ per step, linear convergence $O((1-1/\kappa)^t)$
    \item \textbf{Newton's method}: Second-order, $O(d^3)$ per step (Hessian inversion), quadratic convergence near optimum
    \item \textbf{Why not Newton for DNNs}: (1) Hessian is $d \times d$ where $d$ can be billions, (2) Loss is non-convex so Hessian may not be positive definite, (3) Per-iteration cost is prohibitive
    \item \textbf{Practical approximations}: Adam approximates diagonal Hessian inverse; K-FAC approximates block-diagonal; natural gradient uses Fisher information matrix
    \item \textbf{When to use Newton}: Small-scale convex problems (logistic regression with few features), or as inner loop in constrained optimization
\end{itemize}

\vspace{0.5em}
\textbf{Question 5: How would you debug a model that's not converging?}

\textit{Expected follow-ups:}
\begin{itemize}
    \item What are the most common causes of training failure?
    \item How do you differentiate between optimization and generalization problems?
    \item What monitoring would you set up to diagnose issues early?
\end{itemize}

\textit{Red flags to avoid:}
\begin{itemize}
    \item Immediately suggesting architecture changes without checking basics
    \item Not checking for NaN/Inf values in gradients
    \item Forgetting to verify the data pipeline is correct
\end{itemize}

\textit{Strong answer key points:}
\begin{itemize}
    \item \textbf{Check basics first}: (1) Data pipeline correct? (2) Labels match inputs? (3) Loss function appropriate? (4) Gradients flowing (not NaN/zero)?
    \item \textbf{Learning rate issues}: Loss oscillating $\Rightarrow$ LR too high; loss barely moving $\Rightarrow$ LR too low
    \item \textbf{Gradient issues}: Monitor gradient norms. Exploding $\Rightarrow$ add clipping or reduce LR. Vanishing $\Rightarrow$ check initialization, consider skip connections
    \item \textbf{Overfit a small batch}: If you can't fit 10 examples perfectly, there's a bug or fundamental issue
    \item \textbf{Monitor}: Loss, gradient norms, weight norms, learning rate schedule, validation loss for generalization gap
    \item \textbf{Systematic approach}: Simplify until it works, then add complexity back
\end{itemize}

\end{interviewquestions}

%------------------------------------------------------------------------------
%  CHAPTER SUMMARY
%------------------------------------------------------------------------------
\begin{chaptersummary}
\textbf{What We Covered:}
\begin{itemize}
    \item Vanilla gradient descent suffers on ill-conditioned problems; convergence depends on condition number $\kappa$
    \item Momentum accumulates gradient history, accelerating convergence and dampening oscillations
    \item Nesterov momentum uses look-ahead gradients for earlier correction
    \item Diagnosing momentum issues: effective LR is $\eta/(1-\gamma)$, oscillation patterns, overshoot detection
\end{itemize}

\textbf{Key Formulas:}
\begin{align*}
    &\text{GD Update: } \vect{\theta}_{t+1} = \vect{\theta}_t - \eta \grad L(\vect{\theta}_t) \\
    &\text{Momentum: } \vect{v}_{t+1} = \gamma \vect{v}_t + \grad L(\vect{\theta}_t), \quad \vect{\theta}_{t+1} = \vect{\theta}_t - \eta \vect{v}_{t+1} \\
    &\text{Convergence Rate: } \norm{\vect{\theta}_t - \vect{\theta}^*} \leq \left(\frac{\kappa - 1}{\kappa + 1}\right)^t \norm{\vect{\theta}_0 - \vect{\theta}^*}
\end{align*}

\textbf{When to Use This:}
\begin{itemize}
    \item Core training algorithm for all neural networks---understanding GD is essential
    \item SGD+Momentum for CNNs when generalization is critical
    \item Diagnosing training failures: oscillation suggests LR too high; slow progress suggests LR too low
\end{itemize}

\textbf{Next Up:} Chapter~\ref{ch:advanced-optimizers} covers adaptive optimizers (Adam, AdamW), learning rate schedules, and practical optimizer selection.
\end{chaptersummary}

\end{document}
